\documentclass[10pt,a4paper]{amsart}
\usepackage[margin=19mm]{geometry}
\usepackage{amsmath,amssymb,mathtools}
\usepackage[scr=rsfs,cal=euler]{mathalfa}
\usepackage{xfrac}
\usepackage[unicode,hidelinks,bookmarks=false]{hyperref}
\newcommand{\ie}{i.e.\ }
\newcommand{\cf}{cf.\ }
\newcommand{\resp}{respectively\ }
\newcommand{\Subsection}{\subsection}
\providecommand{\nobreakdash}{\nobreak\hbox{-}}
\setlength{\emergencystretch}{2em}
\multlinegap0pt
\usepackage{bm}
\usepackage{bbm}
\usepackage{dsfont}
\usepackage{xspace}
\usepackage{cancel}
\usepackage{mathtools}
\usepackage{amsthm}
\makeatletter
\@ifundefined{lem}{\newtheorem{lem}{\bf Lemma}}{}
\makeatother
\makeatletter
\expandafter\ifx\csname acknowledgement\endcsname\relax
\newenvironment{acknowledgement}{%
    \titlepage
    \null\vfil
    \@beginparpenalty\@lowpenalty
    \begin{center}%
      \bfseries \ackname
      \@endparpenalty\@M
    \end{center}}%
  {\par\vfil\null\endtitlepage}
\fi
\makeatother
\title[Finite groups of non-prime-power order with exactly four cha]{Finite groups of non-prime-power order with exactly four character codegrees}
\author{Yu Zeng, Mehdi Ghaffarzadeh, Mohsen Ghasemi, Dongfang Yang}
\date{Source: arXiv:2510.08221v1 (CC BY 4.0)}
\begin{document}
\maketitle

\noindent\emph{Source and licence.} Finite groups of non-prime-power order with exactly four character codegrees. Version arXiv:2510.08221v1 (CC BY 4.0), distributed under the Creative Commons Attribution 4.0 International licence, as recorded in the arXiv licence metadata for this exact version and shown on the abstract page https://arxiv.org/abs/2510.08221v1. Full rights notice: licenses/arxiv-2510.08221-CC-BY-4.0.txt.\par\smallskip

\noindent\textit{Editorial scope.} Counted unit: the Lem whose source label is \texttt{lem: nilpotent} in the article (source0.tex lines 2206--2211, source label \texttt{lem: nilpotent}), delivered together with the article's own complete original proof of it (source0.tex lines 2212--2247, one \texttt{proof} environment of 36 lines). No mathematics is written, repaired or re-derived: statement and proof are the article's own text line for line. Required context retained uncounted: lem: direct product, kernel and codegrees (lines 551--564); lem: kernel, abelian, cod (lines 587--596). Every definition, display and label of those lines is retained unchanged. Omitted from this package and not redistributed: 1--550, 565--586, 597--2205, 2248--2548. Nothing inside a retained excerpt is omitted or altered: every retained body is delivered exactly as the article's lines are written, so no omission receipt of any kind accompanies this package. Citations: the retained excerpts carry no citation command at all, so no bibliography entry of the article is transcribed and no reference is redistributed. Reference adaptation: none was needed and none was made; every reference inside the delivered unit resolves inside it, so no unresolved reference or citation is produced. Printed slips kept literal: no printed word, symbol, hypothesis or number of the counted statement or of its proof was altered. Layout and class: the article's own class is \texttt{article} with packages including latexsym, amsfonts, euscript, amsthm, amssymb, stmaryrd, mathrsfs, amsmath, leftidx, mathabx, rotating, xcolor, multirow, booktabs; the wrapper is the lane's pinned \texttt{amsart} editorial wrapper at 10pt on A4, which restates the theorem-like environments the retained text needs and adds the article's own macro definitions that the retained text needs (none), with extra wrapper packages (bm, bbm, dsfont, xspace, cancel, mathtools). The delivered numbering is the unit's own. Assets: no retained span contains a figure environment, a graphics-inclusion command, a PDF-page inclusion, a table or a TikZ picture; no image file is copied into this package, the package declares no asset dependency and no image file is redistributed with it. Licence: version arXiv:2510.08221v1 of this article is distributed under the Creative Commons Attribution 4.0 International licence, as recorded in the arXiv licence metadata for this exact version and shown on the abstract page https://arxiv.org/abs/2510.08221v1. A full rights notice is delivered at licenses/arxiv-2510.08221-CC-BY-4.0.txt. Characters: every retained excerpt is pure ASCII, so the delivered engine, XeLaTeX with its default Latin Modern text font, reports no missing character.
\begin{lem}\label{lem: direct product, kernel and codegrees}
	Let $G$ be the direct product of finite groups $A$ and $B$.
	Then the following hold.
	\begin{description}
		\item[(1)] For $\alpha \in \mathrm{Irr}(A)$ and $\beta \in \mathrm{Irr}(B)$, 
		we have $\ker(\alpha\times \beta)=\ker(\alpha)\times \ker(\beta)$ if and only if $|\mathbf{Z}(\alpha)/\ker(\alpha)|$ and $|\mathbf{Z}(\beta)/\ker(\beta)|$ are coprime. 
		\item[(2)] 
		For $\alpha \in \mathrm{Irr}(A)$ and $\beta \in \mathrm{Irr}(B)$,
		if $(|\mathbf{Z}(\alpha)/\ker(\alpha)|,|\mathbf{Z}(\beta)/\ker(\beta)|)=1$,
		then $\mathrm{cod}(\alpha\times \beta)=\mathrm{cod}(\alpha)\mathrm{cod}(\beta)$.
		\item[(3)] If $|A|$ and $|B|$ are coprime, then $\mathrm{cod}(G)=\{ \mathrm{cod}(\alpha)\mathrm{cod}(\beta):\alpha \in \mathrm{Irr}(A),\beta\in \mathrm{Irr}(B) \}$.
		In particular, $|\mathrm{cod}(G)|=|\mathrm{cod}(A)|\cdot |\mathrm{cod}(B)|$.
	\end{description} 
\end{lem}

\begin{lem}\label{lem: kernel, abelian, cod}
  Let $G$ be a finite group with a normal subgroup
$V = V_{1} \times \dots \times V_{t}$,
where each $V_{i}$ is an abelian minimal normal subgroup of $G$ and
the $V_{i}$ are pairwise non-isomorphic as $G$-modules.
If $\lambda = \lambda_{1} \times \dots \times \lambda_{t}$ with
$\lambda_{i} \in \mathrm{Irr}(V_{i})^{\sharp}$ for every $i$, then
$\ker(\chi) \cap V = 1$ and $|V|$ divides $\mathrm{cod}(\chi)$ for every
$\chi \in \mathrm{Irr}(G | \lambda)$.
\end{lem}

\begin{lem}\label{lem: nilpotent}
  Let $G$ be a solvable group and $N$ the nilpotent residual of $G$.
  Assume that $|\pi(G/N)|>1$.
  Then $|\mathrm{cod}(G)|=4$ if and only if $G$ is a direct product of an elementary abelian $p$-group
  and an elementary abelian $q$-group where $p\neq q$.
\end{lem}

\begin{proof}
  Assume that $G=P\times Q$ where $P$ is a nontrivial elementary abelian $p$-group, $Q$ is 
  a nontrivial elementary abelian $q$-group, and $p\neq q$.
  As $\mathrm{cod}(P)=\{ 1,p \}$ and $\mathrm{cod}(Q)=\{ 1,q \}$,
  we conclude by part (3) of Lemma \ref{lem: direct product, kernel and codegrees} that $\mathrm{cod}(G)=\{ 1,p,q,pq \}$.

  Assume that $|\mathrm{cod}(G)|=4$.
  Since $G/N$ is a nilpotent group with $|\pi(G/N)|>1$,
  it follows by part (3) of Lemma \ref{lem: direct product, kernel and codegrees} that $G/N$ is a direct product of an elementary abelian $p$-group
  and an elementary abelian $q$-group where $p\neq q$.
  Therefore, $N=G'$ and $\mathrm{cod}(G)=\mathrm{cod}(G/N)=\{ 1,p,q,pq \}$.
  
  We next show that $N=1$.
  Let $G$ be a counterexample of minimal possible order such that $N>1$.
 Then $N$ is minimal normal in $G$.
  By Lemma \ref{lem: kernel, abelian, cod}, $|N|\mid \mathrm{cod}(\chi)$ for some $\chi \in \mathrm{Irr}(G)$.
  Therefore, $|N|=p$ or $q$.
  Without loss of generality,
  we may assume that $|N|=p$.
  Then $G=P  \rtimes Q$ where $P\in \mathrm{Syl}_{p}(G)$ and $Q\in \mathrm{Syl}_{q}(G)$.
  Note that $P/N$ is elementary abelian,
  and so $\Phi(P)\leq N$.
  Since $Q$ acts coprimely on $P$, we conclude that $P$ is elementary abelian.
  In fact, otherwise, $\Phi(P)=N$; as $Q$ fixes every element in $P/\Phi(P)$, $[P,Q]=1$, a contradiction.
  So, $G=(N\times P_0) \rtimes Q$ where $P_0\unlhd G$.
  Let $C=\mathbf{C}_{Q}(N)$ and $Q_0$ a complement of $C$ in $Q$.
  Then $NQ_0$ is a Frobenius group with complement $Q_0\cong \mathsf{C}_{q}$ and kernel $N\cong \mathsf{C}_{p}$.
  Also, $G= NQ_0 \times (P_0\times C)$.
  Let $\theta$ be a nonlinear character in $\mathrm{Irr}(NQ_0)$ and $\beta\in \mathrm{Irr}(P_0)^\sharp$.
  Then $\chi=\theta \times (\beta\times 1_C) \in \mathrm{Irr}(G)$.
  Since $NQ_0$ is a Frobenius group and $\ker(\theta)=1$,
  $\mathbf{Z}(\theta)=\mathbf{Z}(NQ_0)=1$.
  So, it follows by part (2) of Lemma \ref{lem: direct product, kernel and codegrees} that 
  $\mathrm{cod}(\chi)=\mathrm{cod}(\theta)\mathrm{cod}(\beta\times 1_C)=\mathrm{cod}(\theta)\mathrm{cod}(\beta)\mathrm{cod}(1_C)=p^{2}$,
  a contradiction.
\end{proof}

\end{document}
